\documentclass[a4paper,12pt]{article}
\usepackage{docsTemplate}
\usepackage{eurosym}

\setTitle{Analisi dei requisiti}
\setAuthors{---}
\setVerificators{Nenad Radulovic}
\setApprovation{Bianca Zaghetto}
\setVersion{0.1}
\setType{Esterno}
\setDestination{Tullio Vardanega \\ & Riccardo Cardin \\ & Team Three Technologies}

\begin{document}
    \include{docTitlepage}

    \section*{Registro delle modifiche} {
        \begin{table}[h!]
            \rowcolors{2}{lightgray}{white}
                \begin{tabularx}{\textwidth}{|l|l|l|l|L|L|}
                \hline
                \textbf{Vers.} & \textbf{Data} & \textbf{Autore} & \textbf{Ruolo} & \textbf{Verifica} & \textbf{Oggetto} \\
                \hline
                \textbf{0.1} & 7/11/25 & Sara Gioia Fichera & Analista &  & Struttura iniziale del documento \\
                \hline
            \end{tabularx}
        \end{table}
    }
    \newpage
    
    \tableofcontents
    \newpage
    
    \section{Introduzione} {
        \subsection{Scopo del documento}
        
            Il presente documento ha lo scopo di descrivere in modo esaustivo i casi d’uso e i requisiti del progetto DIPReader.
            Tale analisi è il risultato di un attento studio del capitolato, integrato con quanto concordato durante gli incontri svolti con Sanmarco Informatica, il proponente.

        \subsection{Glossario}
            Si rimanda al documento Glossario trovabile nella cartella Documenti/1-RTB. I vocabili presenti nel glossario sono stati segnalati con il seguente stile: colore. 
        \subsection{Riferimenti}
            \subsubsection{Riferimenti normativi}
                \begin{itemize}
                    \item Norme di progetto
                    \item Capitolato
                    \begin{itemize}
                        \item \href{https://www.math.unipd.it/~tullio/IS-1/2025/Progetto/C3.pdf}{\textcolor{blue} {https://www.math.unipd.it/~tullio/IS-1/2025/Progetto/C3.pdf}}
                        \item[] (Ultimo accesso: 2025/11/07)
                    \end{itemize}
                    \item Regolamento del Progetto Didattico
                    \begin{itemize}
                        \item \href{https://www.math.unipd.it/~tullio/IS-1/2025/Dispense/PD1.pdf}{\textcolor{blue} {https://www.math.unipd.it/~tullio/IS-1/2025/Dispense/PD1.pdf}}
                        \item[] (Ultimo accesso: 2025/11/07)
                    \end{itemize}
                \end{itemize}
            \subsubsection{Riferimenti informativi}
                \begin{itemize}
                    \item Slide del corso: T4 - Gestione di progetto
                    \begin{itemize}
                        \item \href{https://www.math.unipd.it/~tullio/IS-1/2025/Dispense/T04.pdf}{\textcolor{blue} {https://www.math.unipd.it/~tullio/IS-1/2025/Dispense/T04.pdf}}
                        \item[] (Ultimo accesso: 2025/11/07)
                    \end{itemize}
                    \item Slide del corso: T5 - Analisi dei requisiti
                    \begin{itemize}
                        \item \href{https://www.math.unipd.it/~tullio/IS-1/2025/Dispense/T05.pdf}{\textcolor{blue} {https://www.math.unipd.it/~tullio/IS-1/2025/Dispense/T05.pdf}}
                        \item[] (Ultimo accesso: 2025/11/07)
                    \end{itemize}
                    \item Slide del corso: P2 - Diagrammi dei Casi d'Uso
                        \begin{itemize}
                            \item \href{https://www.math.unipd.it/~rcardin/swea/2022/Diagrammi\%20Use\%20Case.pdf}{\textcolor{blue}{https://www.math.unipd.it/~rcardin/swea/2022/Diagrammi\%20Use\%20Case.pdf}}
                            \item[] (Ultimo accesso: 2025/11/07)
                        \end{itemize}
                    \item Glossario
                    \item Altra documentazione...
                \end{itemize}
    }

    \newpage
    \section{Descrizione del prodotto} {
        \subsection{Presentazione del prodotto}
            DIPReader ha lo scopo di supportare la gestione e la conservazione digitale dei documenti, consentendo una transizione strutturata da archivi cartacei a soluzioni cloud avanzate.
            \\Il sistema centralizza i documenti e ne assicura autenticità, integrità, leggibilità, reperibilità e sicurezza, offrendo funzionalità di acquisizione, estrazione dei dati e consultazione. Inoltre, permette di richiedere la disponibilità offline dei documenti, garantendo sempre la conformità ai requisiti normativi. 
        \subsection{Funzioni del prodotto}
        Le funzioni del prodotto sono: 
        \begin{itemize}
            \item Autenticazione
            \item Rappresentazione semplificata del contenuto 
            \item Ricerca dei documenti tramite metadati
            \item Visualizzazione in anteprima di alcuni formati
            \item Salvare documenti nel computer dell'Utente
            \item Funzionamento offline
        \end{itemize}
        \subsection{Caratteristiche dell'utente}

    }

    \section{Casi d'uso}
        \subsection{Scopo}
            Lo scopo di questa sezione è spiegare singolarmente ciascun Caso d'uso individuato.  

        \subsection{Attori}
            \begin{itemize}
                \item Utente:
            \end{itemize}
        
        \subsection{Casi d'uso}
            \subsubsection{UC1}
            \subsubsection{UC2}

    \section{Requisiti}
            In questa sezione vengono illustrati i requisiti che stanno alla base di DIPReader. Essi sono classificati nella forma: XX.XX.XX. . 
        \subsection{Requisiti funzionali}
        \subsection{Requisiti di qualità}
        \subsection{Requisiti di vincolo}


\end{document}